\documentclass[../../main.tex]{subfiles}
\begin{document}

平面上の2点 $P$，$Q$ の距離を $d(P,Q)$ と表すことにする．
平面上に点 $O$ を中心とする一辺の長さが1000の正三角形 $\triangle A_1A_2A_3$ がある．
$\triangle A_1A_2A_3$ の内部に3点 $B_1$，$B_2$，$B_3$ を，$d(A_n,B_n)=1$ $(n=1,2,3)$ となるようにとる．
また，
\begin{align*}
  \vec{a}_1=\overrightarrow{A_1A_2}, \hspace{1\zw} \vec{a}_2=\overrightarrow{A_2A_3}, \hspace{1\zw} \vec{a}_3=\overrightarrow{A_3A_1} \\
  \vec{e}_1=\overrightarrow{A_1B_1}, \hspace{1\zw} \vec{e}_2=\overrightarrow{A_2B_2}, \hspace{1\zw} \vec{e}_3=\overrightarrow{A_3B_3}
\end{align*}
とおく．
$n=1,2,3$ のそれぞれに対して，時刻0に $A_n$ を出発し，
$\vec{e}_n$ の向きに速さ1で直進する点を考え，時刻 $t$ におけるその位置を $P_n(t)$ と表すことにする．

\begin{figure}[htbp]
  \centering
  \begin{tikzpicture}[scale=1.3]
    \coordinate (O) at (0,0);
    \coordinate (A1) at (90:2.5);
    \coordinate (A2) at (210:2.5);
    \coordinate (A3) at (330:2.5);

    \draw[thick] (A1) -- (A2) -- (A3) -- cycle;

    \node[circle, fill, inner sep=1.5pt] at (O) {}; \node[below right] at (O) {$O$};

    \coordinate (P1) at ($(A1) + (-75:1.2)$);
    \draw[thick] (A1) -- (P1);
    \node[circle, fill, inner sep=1.5pt] at (P1) {}; \node[below] at (P1) {$P_1(t)$};

    \coordinate (P2) at ($(A2) + (25:1.2)$);
    \draw[thick] (A2) -- (P2);
    \node[circle, fill, inner sep=1.5pt] at (P2) {}; \node[above] at (P2) {$P_2(t)$};

    \coordinate (P3) at ($(A3) + (145:1.2)$);
    \draw[thick] (A3) -- (P3);
    \node[circle, fill, inner sep=1.5pt] at (P3) {}; \node[above] at (P3) {$P_3(t)$};

    \draw (A1) +(-120:0.5) arc (-120:-75:0.5);
    \node at ($(A1)+(-100:0.7)$) {$\theta_1$};

    \draw (A2) +(0:0.5) arc (0:25:0.5);
    \node at ($(A2)+(12:0.7)$) {$\theta_2$};

    \draw (A3) +(120:0.5) arc (120:145:0.5);
    \node at ($(A3)+(132:0.7)$) {$\theta_3$};

    \node[above] at (A1) {$A_1$};
    \node[below left] at (A2) {$A_2$};
    \node[below right] at (A3) {$A_3$};
  \end{tikzpicture}
\end{figure}

\begin{enumerate}
  \item ある時刻 $t$ で $d(P_1(t),P_2(t))\leqq1$ が成立した．
    ベクトル $\vec{e}_1-\vec{e}_2$ と，ベクトル $\vec{a}_1$ とのなす角度を $\theta$ とおく．
    このとき $\displaystyle|\sin\theta|\leqq\dfrac{1}{1000}$ となることを示せ．

  \item 角度 $\theta_1$，$\theta_2$，$\theta_3$ を
    $\theta_1=\angle B_1A_1A_2$，$\theta_2=\angle B_2A_2A_3$，$\theta_3=\angle B_3A_3A_1$ によって定義する．
    $\alpha$ を $\displaystyle0<\alpha<\dfrac{\pi}{2}$ かつ $\displaystyle\sin\alpha=\dfrac{1}{1000}$ をみたす実数とする．
    (1) と同じ仮定のもとで，$\theta_1+\theta_2$ の値のとる範囲を $\alpha$ を用いて表せ．

  \item 時刻 $t_1$，$t_2$，$t_3$ のそれぞれにおいて，次が成立した．
    \begin{align*}
      d(P_2(t_1),P_3(t_1))\leqq1, \hspace{1\zw} d(P_3(t_2),P_1(t_2))\leqq1, \hspace{1\zw} d(P_1(t_3),P_2(t_3))\leqq1
    \end{align*}
    このとき，時刻 $\displaystyle T=\dfrac{1000}{\sqrt3}$ において同時に
    \begin{align*}
      d(P_1(T),O)\leqq3, \hspace{1\zw} d(P_2(T),O)\leqq3, \hspace{1\zw} d(P_3(T),O)\leqq3
    \end{align*}
    が成立することを示せ．
\end{enumerate}

\end{document}
